본 연구는 문항 기반 검사를 중심으로 유효숫자 적용 양상을 분석하였다는 점에서
실제 실험 및 탐구 상황을 완전히 반영하지 못한 한계가 있다.
%
학생들이 도구를 직접 조작하며 판단하는 과정이나, 탐구 보고서, 기록물에서 나타나는
자연스러운 자릿수 선택 행동은 포함되지 않았기에,
측정 맥락에서 드러나는 실제 적용 수준을 모두 확인하기 어렵다.
%
또한 본 연구에서 제시한 측정 도구와 상황이 제한적이어서,
다양한 도구 사용 경험이나 탐구 맥락에 따라 변하는 판단 차이를 포괄적으로 파악하기에는 제약이 있다.
%
이와 같은 점은 본 연구가 문항 기반 검사로 파악할 수 있는 범위에 한정된다는 특성을 보여주며,
실제 탐구 맥락에서의 적용 수준을 정교하게 파악하기 위해서는 추가적인 자료 수집이 필요하다.